\section{Evaluación, confiabilidad y seguridad: de la demo al sistema en producción}

Esta entrevista no deja el problema del Agent en la pregunta «¿se puede hacer una demo impresionante?». Su Yu (苏煜) reconduce varias veces la discusión hacia reliability, speed, cost effectiveness y deployment. Una demo de agent puede mostrar que «sabe hacerlo» mediante un entorno, unos ejemplos y unas indicaciones humanas diseñados con cuidado; pero un sistema en producción tiene que responder otras preguntas: qué tasa de fallos tiene ante tareas desconocidas, si los fallos son recuperables, si el costo es controlable, si los límites de permisos están claros y cómo se acotan las fugas de datos o las operaciones erróneas.

\begin{importantbox}{Cuatro umbrales para pasar de la demo a producción}
El primero es la tasa de éxito de la tarea, no solo que cada acción individual sea correcta; el segundo es la recuperabilidad, es decir, si tras un fallo se puede localizar el error y continuar; el tercero es el costo, si los token, las llamadas a herramientas y el respaldo humano de las tareas de largo alcance resultan económicos; el cuarto es la seguridad: cuando el agent puede operar sobre entornos reales, los permisos, la auditoría y el peor escenario posible tienen que incorporarse desde el diseño.
\end{importantbox}

\subsection{Por qué el benchmark no es suficiente}

Los primeros benchmark de web agent o de computer-use permitieron a los investigadores comparar sistemas, pero con frecuencia no cubren por completo un despliegue real. En el entorno real de una empresa hay sesiones iniciadas, permisos, datos históricos, procesos organizacionales, sistemas heredados, requisitos de auditoría e instrucciones incompletas de los usuarios. Que un agent tenga un buen desempeño en un benchmark estándar no significa que pueda operar de forma segura en un banco, en el área legal, en el sector salud, en una cadena de herramientas de investigación y desarrollo o en los sistemas internos de una empresa.

\noindent\begin{tabularx}{\textwidth}{p{0.20\textwidth}p{0.34\textwidth}X}
\toprule
Dimensión de evaluación & Qué hay que medir & Por qué importa \\
\midrule
Task success & Si la tarea de largo alcance se completa al final & Evita optimizar solo un clic aislado o una respuesta única. \\
Recovery & Si tras un error puede autodiagnosticarse, revertir y reintentar & Determina si el sistema puede operar sin supervisión. \\
Cost & Token, llamadas a herramientas, tiempo de espera, intervención humana & Determina si el agent puede desplegarse a escala. \\
Safety & Límites de permisos, acciones peligrosas, privacidad y auditoría & Condición necesaria para pasar de la demo a la operación real del negocio. \\
Adaptation & Si puede aprender de entornos nuevos y de retroalimentación nueva & Vincula el continual learning con el expert agent. \\
\bottomrule
\end{tabularx}

\subsection{La diferencia entre Security y Safety}

En la entrevista aparece una distinción que se pasa por alto con facilidad: muchos problemas de safety son, en el fondo, problemas de capacidad, porque el agent, como un principiante, no entiende el entorno y comete errores elementales; en cambio, la security se orienta más al worst-case scenario y requiere métodos específicos. Es decir, aumentar la capacidad puede reducir muchos errores por descuido, pero no sustituye el control de permisos, los entornos aislados, la auditoría ni los ejercicios de equipo rojo.

\begin{warningbox}{Mejorar la capacidad no es diseñar la seguridad}
Un agent más inteligente quizá haga menos clics equivocados, pero también puede ser más capaz de eludir las restricciones. Un sistema en producción debe diseñar por separado la mejora de capacidad y los mecanismos de seguridad: la capacidad se encarga de completar la tarea; la seguridad, de acotar el espacio de la tarea, registrar el comportamiento y bloquear las acciones inaceptables.
\end{warningbox}


\subsection{Lista práctica: qué preguntar antes de construir un sistema de Agent}

Si el contenido de esta entrevista se convierte en una checklist de ingeniería, hay que plantear al menos ocho preguntas. Primera: cuáles son los límites del Agent, quién puede asignarle objetivos y quién puede detenerlo. Segunda: cuál es el entorno al que se enfrenta y cómo se observa el estado de ese entorno. Tercera: qué acciones puede realizar y cuáles deben prohibirse o requerir confirmación humana. Cuarta: cómo se determina automáticamente el éxito de la tarea y qué tareas requieren aceptación humana. Quinta: cómo se recupera tras un fallo y si puede revertir cambios o generar registros de auditoría. Sexta: cómo acumula experiencia y si esa experiencia puede contaminar tareas posteriores. Séptima: si su costo es predecible. Octava: desde qué punto de entrada del producto lo invocarán los usuarios de forma natural.

\begin{importantbox}{Criterio de ingeniería}
Que un sistema de Agent pueda ponerse en producción no depende de que complete una tarea asombrosa en una demostración, sino de que las ocho preguntas anteriores puedan responderse de forma estable. Esta lista también sirve como marco reutilizable para leer otras entrevistas sobre Agent.
\end{importantbox}

\subsection{Resumen de la sección}

Llevar un Agent a producción no es solo un problema del modelo. La evaluación, la recuperación, el costo, la seguridad y la adaptación al entorno determinan en conjunto si puede pasar de la demo a un sistema real. Por eso el continual learning y el world model se vuelven el eje central: en última instancia deben servir a la confiabilidad, no solo a demostraciones más vistosas.
